% ============================================================
%  HCMUT-DEE Thesis Kit v2.0.3 — Phụ lục mẫu
%  Copyright (c) 2026 Nguyen Trong Thang & TS. Nguyen Phuc Khai
%  License: MIT
% ============================================================

\chapter{Phụ lục mẫu}
\label{app:sample}

% ═══════════════════════════════════════════════════════════════
%  GHI CHU: Day la phu luc mau. Hay thay the bang noi dung
%  phu luc cua do an ban. Vi du: ma nguon, du lieu, bang phu.
% ═══════════════════════════════════════════════════════════════

\section{Ví dụ mã nguồn Python}

\begin{lstlisting}[style=pythonstyle, caption={Hàm tính tổn thất công suất trên lưới điện}, label={lst:app-python}]
import numpy as np

def backward_forward_sweep(R, X, P_load, Q_load, V_base=12.66):
    """
    Giai luu luong cong suat bang phuong phap
    Backward/Forward Sweep cho luoi hinh tia.
    """
    n_bus = len(P_load)
    V = np.ones(n_bus) * V_base  # Khoi tao dien ap

    for iteration in range(100):
        # Backward sweep: tinh dong nhanh
        I_branch = np.zeros(n_bus - 1, dtype=complex)
        for k in range(n_bus - 2, -1, -1):
            S = complex(P_load[k+1], Q_load[k+1])
            I_branch[k] = np.conj(S / V[k+1])

        # Forward sweep: cap nhat dien ap
        for k in range(n_bus - 1):
            Z = complex(R[k], X[k])
            V[k+1] = V[k] - I_branch[k] * Z

        # Kiem tra hoi tu
        if np.max(np.abs(V - V_old)) < 1e-6:
            break
        V_old = V.copy()

    return V
\end{lstlisting}

\section{Ví dụ bảng dữ liệu}

\begin{table}[H]
\centering
\caption{Thông số hệ thống IEEE 33-bus (trích một phần)}
\label{tab:app-ieee33}
\begin{tabular}{cccccc}
\toprule
\textbf{Nhánh} & \textbf{Từ} & \textbf{Đến} & \textbf{R ($\Omega$)} & \textbf{X ($\Omega$)} & \textbf{P$_{load}$ (kW)} \\
\midrule
1  & 1 & 2  & 0.0922 & 0.0477 & 100 \\
2  & 2 & 3  & 0.4930 & 0.2511 & 90  \\
3  & 3 & 4  & 0.3661 & 0.1864 & 120 \\
4  & 4 & 5  & 0.3811 & 0.1941 & 60  \\
5  & 5 & 6  & 0.8190 & 0.7070 & 60  \\
\bottomrule
\end{tabular}
\end{table}

% === CHEN FILE PDF BEN NGOAI (bo comment neu can) ===
% \section{Datasheet thiet bi}
% \includepdf[pages=-]{path/to/datasheet.pdf}
